\section{Lessons and Threats to Validity}\label{sec:threats}

% ICDE: summary of findings and recommendations
{\footnotesize
\begin{xltabular}{\textwidth}{@{}l >{\raggedright\arraybackslash}X >{\raggedright\arraybackslash}p{0.27\textwidth} >{\raggedright\arraybackslash}p{0.13\textwidth}@{}}
\caption{Summary of findings and recommendations.}\label{tab:findings}\\
\toprule
& Finding (evidence) & Recommendation & Audience\\
\midrule
\endfirsthead
\multicolumn{4}{@{}l}{\emph{Table~\ref{tab:findings}, continued}}\\
\toprule
& Finding (evidence) & Recommendation & Audience\\
\midrule
\endhead
\bottomrule
\endfoot
F1 & Under one SQLite certifier, released counterexamples lose about half of their pairs at the DDL tier and 82--92\,\% at the format tier. SpotIt, re-run under our budget, certifies \V{srLonef} pairs at $\mathcal A_{1f}$, against CaliberGraph's \V{nCGLonef}. The ranking of tools changes with the tier (Sec.~\ref{sec:main}). & Compare disprovers only under one certifier, tier and budget. Re-certify released witnesses on the benchmark's engine. & tool authors\\\addlinespace
F2 & By an LLM attribution, \V{attrImplMin}--\V{attrImplMax}\,\% of every tool's schema-tier witnesses rely on implausible data. Two humans shown facts of the data agreed (\V{hvTwoImpConsOne}/\V{hvTwoConsOne}). However, three further humans agreed with each other barely above chance (Fleiss' $\kappa_2=\V{hvThreeFleissBin}$), and their majority called \V{hvThreeImpOne}/\V{hvThreeMajOne} implausible. Humans attribute more witnesses to prediction errors than the AI does. A repair test removes \V{xVanPct}\,\% of the differences (Sec.~\ref{sec:validity}). & Report witnesses with an attributed sample, the data facts the judgment rests on and the annotators' agreement. Never report bare \emph{non-equivalent} counts. & tool authors, maintainers\\\addlinespace
F3 & A data profile of the benchmark takes minutes to mine. It cuts implausible witnesses from \V{pCGImpl}\,\% to \V{pCGRImpl}\,\% by the AI attribution; the human majority rates \V{hvThreeImpOne}/\V{hvThreeMajOne} vs.\ \V{hvThreeImpTwo}/\V{hvThreeMajTwo} as implausible. SpotIt+'s counterexamples (LLM-validated column constraints) are \V{pSIPImpl}\,\% implausible. Re-certified counts are insensitive to the profile's thresholds (Sec.~\ref{sec:rresults}). & Publish the data profile with the benchmark. Judge witnesses at $\mathcal A_R$, with a presupposition constraint for entities named in the question. & maintainers\\\addlinespace
F4 & At $\mathcal A_R$, counting witnesses lowers EX by \V{ptsCountRMin}--\V{ptsCountRMax} points per system, against 15--20 at $\mathcal A_{1f}$. Weighting by attributed error rates gives \V{ptsErrRMin}--\V{ptsErrRMax} points with the AI labels or \V{ptsHumThreeRMin}--\V{ptsHumThreeRMax} with the round-3 human majority (Sec.~\ref{sec:rankings}). & Never subtract witness counts from EX. Use meaningful witnesses to repair gold queries and questions. & maintainers, system builders\\\addlinespace
F5 & Join-graph completion and real-row retrieval matter, but three synthetic components do not. A data-grounded audit of one system's accepted test predictions costs under one core-hour at 5\,s per pair, plus attribution (Sec.~\ref{sec:ablation}). & Build disprovers on the join graph and real data. Certify and attribute every leaderboard submission. & tool authors, maintainers\\\addlinespace
F6 & On Spider, schema-tier witnesses again fail the profile, but CaliberGraph-R's lead does not replicate and fuzzing-R's witnesses are more meaningful. On 13 unseen BIRD systems, the lead in counts replicates (Sec.~\ref{sec:replication}). & Mine profiles from instances large enough to expose keys and dependencies, or add them by hand. & maintainers, tool authors\\
\end{xltabular}
}

\textbf{Lessons.} Table~\ref{tab:findings} maps each finding to a recommendation. The common thread is that a counterexample is a database and should be judged as data: certified on the benchmark's engine with the original query strings, checked against the benchmark's own profile, and released so that others can attribute it. Meaningful witnesses are feedback on the benchmark as well as on the systems. By the AI attribution, 58--64\,\% of the sampled data-grounded witnesses expose latent gold-query defects or ambiguous questions, against 11--22\,\% that expose prediction errors. Human annotators see more prediction errors (\V{humThreeErrPctR}\,\% by the round-3 majority), but they agree too little with each other for a reliable estimate. No correction of system scores can fix gold defects.

\textbf{Disproving is not proving.} A missing witness says nothing. Let \emph{semantic accuracy} at tier $\mathcal A$ be the share of predictions that are equivalent to the gold query under $(\mathcal C,\mathcal A)$. Accuracy corrected with witnesses at $\mathcal A$ is an upper bound on it, since every removed prediction is certified non-equivalent, while predictions with undetected differences remain counted. It is not a bound on the share of predictions that answer the question correctly, because witnesses that expose gold defects or ambiguities also remove predictions that may be right. Error-weighted corrections are lower bounds on the overstatement if the labels are unbiased. Bounded verifiers can also report the absence of counterexamples up to a bound, which our tools cannot do.

\textbf{What \emph{implausible} means.} The profile is mined from one instance, so it can be both too strict and too weak. It is too strict because a plausible but unobserved value is admitted only through range and shape rules. It is too weak because annotators still found conditional constraints that it does not express, such as a badge awarded once while others repeat; conditional-dependency and denial-constraint discovery could capture such constraints~\cite{fan2011cfd,chu2013dc}. \emph{Implausible} combines absence from the instance with a judgment about the domain, and annotators disagree on this judgment (Section~\ref{sec:validity}). A live football database may well contain unplayed fixtures without scores (Fig.~\ref{fig:example}). Conversely, a prediction that is right only because a column happens to hold no \texttt{NULL} is fragile under updates. An implausible witness therefore says that the benchmark, not the domain, cannot separate the queries. The same definition makes H7 close to circular: a judge who can query $D_0$ will rarely call a witness built from real rows of $D_0$ implausible. H7 thus checks that $\mathcal A_R$ works as constructed. The evidence that does not depend on the judge comes from mechanical outcomes: profile violations and repair survival (Section~\ref{sec:rresults}). Sub-databases can also omit the entity a question names (Section~\ref{sec:rresults}). Some differences rely on behavior that SQL leaves unspecified, such as which tied row \texttt{LIMIT} keeps. Because BIRD's EX is computed with SQLite, we count them, and attribution classifies most of them as ambiguities.

\textbf{AI-performed annotation.} All attribution and root-cause labels were produced by the AI model that also implemented the tools, in two independent passes plus adjudication. Agreement between the passes measures consistency, not correctness. The first author's earlier checks of two samples of 30 witnesses (agreement 26--28 of 30) were not blind: the AI verdicts were printed on the same sheet. A second model reproduces the labels (Section~\ref{sec:validity}), but both models share a vendor and the same data access. The human annotators of rounds 1--2, two data analysts familiar with SQL who volunteered, labeled independently. The pre-registered blind check failed in its first round, in which they had no access to $D_0$. In a second, separately pre-registered round with facts from $D_0$, they agreed with each other and reproduced the difference between the tiers, but they agreed with the AI only moderately. The two round-2 sheets contained a few manual edits outside the label fields: end-of-file markers, two altered witness cells in one item and a shortened fact line in another. The authors first attributed these edits to an editing tool. Excluding the two items changes no outcome. In round~3, three further unpaid analysts agreed barely above chance, and the process record of this round has the inconsistencies listed in Section~\ref{sec:validity}. The authors state that all three rounds were annotated by hand, and the round-3 annotators reported working without AI tools. All sheets are released as received, and the AI's estimated shares carry this uncertainty. Blinding to the tier is imperfect, since synthetic values such as \texttt{\textquotesingle{}x1\textquotesingle{}} and \texttt{-1} are recognizable. All annotations, reasons, witnesses and all human sheets are released.

\textbf{Baselines.} SpotIt was re-run from its released code under our budget (Section~\ref{sec:main}). Its bounded solver often exhausts five seconds, so a larger budget would favor it. Its released claims, obtained with larger budgets, are reported as well. SpotIt+ is used through its released outputs, since its constraint mining needs an LLM API that is unavailable in our environment. ParSEval-R certifies nothing at $\mathcal A_R$ and is not a fair representative of symbolic generation there. Unlike every other tool, the frozen first-stage fuzzing baseline accepted the first differing database without consulting the certifier. We re-ran exactly the affected pairs (132 distinct EX-correct, 259 EX-incorrect) with certifier-in-the-loop acceptance and the same seeds, and we report the corrected counts. The attribution sample of this baseline was drawn before the correction.

\textbf{Deviations from the pre-registrations.} The six documented deviations include adding the first-stage attribution after the first test run, adding the repair test post hoc, revising the data-profile tier on development data before the second freeze, and designing the second stage after seeing first-stage test attributions. Appendix~\ref{app:deviations} lists all deviations, and the artifact repository lists every amendment with its timestamp.

\textbf{Generalization.} All results concern SQLite, BIRD and Spider~1.0. The systems answer the same questions, and our intervals account for this by resampling questions. The tier was conceived after first-stage test attributions had been seen. Its only test on databases that no design decision had seen, Spider, did not replicate CaliberGraph-R's lead; the lead replicated only on new systems over the BIRD databases. Spider also shows that a profile mined from small instances is less informative. Larger enterprise schemas and other engines remain to be studied.
